\DocumentMetadata{
  pdfversion=2.0,
  pdfstandard=ua-2,
  lang=en-US,
  tagging=on,
  tagging-setup={math/setup=mathml-SE,tabsorder=structure}
}
\documentclass[11pt,letterpaper]{article}
\usepackage[margin=0.68in,includefoot]{geometry}
\usepackage{newcomputermodern}
\usepackage{amsmath,amscd,mathtools}
\providecommand{\square}{\mdlgwhtsquare}
\usepackage[hidelinks]{hyperref}
\hypersetup{
  pdftitle={Algebra PhD Exam, May 2021},
  pdfauthor={University of Florida Department of Mathematics},
  pdfsubject={Graduate mathematics examination},
  pdfdisplaydoctitle=true,
  bookmarksopen=true
}
\setcounter{secnumdepth}{0}
\setlength{\parindent}{0pt}
\setlength{\parskip}{0.28em}
\setlength{\emergencystretch}{3em}
\setlength{\leftmargini}{2.15em}
\setlength{\leftmarginii}{2.65em}
\setlength{\leftmarginiii}{3em}
\renewcommand{\labelenumi}{\textbf{\theenumi.}}
\pagestyle{empty}
\raggedbottom
\allowdisplaybreaks
\newenvironment{examproblems}[1]{%
  \begin{enumerate}%
  \setcounter{enumi}{#1}%
  \setlength{\topsep}{0.35em}%
  \setlength{\partopsep}{0pt}%
  \setlength{\itemsep}{0.48em}%
  \setlength{\parsep}{0.18em}%
}{\end{enumerate}}
\newenvironment{examsubparts}{%
  \begin{enumerate}%
  \setlength{\topsep}{0.18em}%
  \setlength{\partopsep}{0pt}%
  \setlength{\itemsep}{0.16em}%
  \setlength{\parsep}{0.1em}%
}{\end{enumerate}}
\newenvironment{examinstructions}{%
  \par\begingroup\itshape
}{%
  \par\endgroup\vspace{0.18em}
}
\makeatletter
\renewcommand\section{\@startsection{section}{1}{0pt}%
  {-1.25ex plus -.25ex minus -.1ex}%
  {0.55ex plus .1ex}%
  {\normalfont\large\bfseries}}
\renewcommand{\@maketitle}{%
  \newpage\null\vskip -1.2em%
  {\footnotesize\bfseries
    \href{https://www.ufl.edu/}{University of Florida}, \href{https://math.ufl.edu/}{Department of Mathematics}\par}%
  \vskip 0.18em%
  \tagstructbegin{tag=Title}%
    {\LARGE\bfseries\href{https://gma.math.ufl.edu/exams/algebra-phd/}{Algebra PhD Exam}, May 2021\par}%
  \tagstructend%
  \vskip 0.35em%
  \tagmcbegin{artifact}%
    \hrule height 1pt%
  \tagmcend%
  \vskip 0.55em%
}
\makeatother
\title{Algebra PhD Exam, May 2021}
\author{}
\date{}
\begin{document}
\maketitle
\thispagestyle{empty}
\begin{examinstructions}
Answer 7 questions from among the ones below. If you answer more than that, only the first seven will be graded.

\par
Write your answers clearly in complete English sentences. You may quote results (within reason) as long as you state them clearly.
\end{examinstructions}
\begin{examproblems}{0}
\item
Suppose that~\(p\) is an odd prime and \(K/\mathbb{Q}\) is the extension of~\(\mathbb{Q}\) obtained by adjoining a primitive~\(p\)th root of~\(1\) in~\(\mathbb{C}\).
\begin{examsubparts}
\item[\textnormal{(i)}]
Show that~\(K\) contains a unique quadratic extension~\(F\) of~\(\mathbb{Q}\).
\item[\textnormal{(ii)}]
What is \(\operatorname{Gal}(K/F)\)?
\end{examsubparts}
\item
Let~\(G\) be a group, and recall that a~\(G\)-set is a set~\(S\) with a left~\(G\)-action \(G \times S \to S\). Write the action as \((g,x) \mapsto g.x\). A morphism of~\(G\)-sets \(f:S\to T\) is a map \(f:S\to T\) such that \(f(g.x)=g.f(x)\) for all \(g\in G\) and \(x\in S\).
\begin{examsubparts}
\item[\textnormal{(i)}]
Show that~\(G\)-sets and morphisms of~\(G\)-sets form a category \(\mathbf{G\text{-}Set}\).
\item[\textnormal{(ii)}]
Show that there is a functor \(H:\mathbf{G\text{-}Set}\to\mathbf{Set}\) which assigns to any~\(G\)-set the underlying set.
\item[\textnormal{(iii)}]
Show that there is a functor \(F:\mathbf{Set}\to\mathbf{G\text{-}Set}\) such that \(F(S)\) is a free~\(G\)-set on~\(S\) (endow the set \(G\times S\) with an appropriate~\(G\)-action).
\end{examsubparts}
\item
Suppose~\(A\) is a ring with identity and~\(I\) is a set.
\begin{examsubparts}
\item[\textnormal{(i)}]
Show that if \(\{P_\alpha\}_{\alpha\in I}\) is a set of left~\(A\)-modules then \(\bigoplus_\alpha P_\alpha\) is projective if and only if each \(P_\alpha\) is projective.
\item[\textnormal{(ii)}]
Show that if \(\{M_\alpha\}_{\alpha\in I}\) is a set of left~\(A\)-modules then \(\prod_{\alpha\in I}M_\alpha\) is injective if and only if \(M_\alpha\) is injective.
\end{examsubparts}
\item
Suppose~\(\mathcal{C}\) is a category and~\(X\) and~\(Y\) are objects of~\(\mathcal{C}\).
\begin{examsubparts}
\item[\textnormal{(i)}]
Define what it means for an object~\(Z\) of~\(\mathcal{C}\) to be a coproduct of~\(X\) and~\(Y\).
\item[\textnormal{(ii)}]
Show that any two groups have a coproduct in the category of groups (if~\(G\) and~\(H\) are groups, construct the coproduct of~\(G\) and~\(H\) as a suitable quotient of the free group on the disjoint union \(G\amalg H\)).
\end{examsubparts}
\item
Suppose~\(K\) is a field, \(K^{alg}\) is an algebraic closure of~\(K\) and \(L/K\) is a finite separable extension. Show that \(L\otimes_K K^{alg}\) is a direct sum of fields. Hint: if \(\alpha\in L\) has minimal polynomial \(f\in K[X]\), observe that there is an exact sequence 
\[
0\to K[X]f\to K[X]\to K(\alpha)\to 0.
\]
\item
Suppose~\(A\) is a Dedekind ring and \(S\subseteq A\) is a multiplicative system. Show that \(S^{-1}A\) is a Dedekind ring.
\item
Suppose~\(A\) is a commutative ring with identity and \(S\subseteq A\) is a multiplicative system.
\begin{examsubparts}
\item[\textnormal{(i)}]
Show that for any~\(A\)-module~\(M\), \(S^{-1}M\simeq S^{-1}A\otimes_A M\).
\item[\textnormal{(ii)}]
Show that \(S^{-1}A\) is a flat~\(A\)-module.
\end{examsubparts}
\item
Prove the Hilbert basis theorem: if~\(A\) is a commutative Noetherian ring then so is the polynomial ring \(A[X]\).
\item
Suppose~\(A\) is an integral domain with fraction field~\(K\).
\begin{examsubparts}
\item[\textnormal{(i)}]
Define what it means for an~\(A\)-submodule \(I\subseteq K\) to be an invertible fractional ideal.
\item[\textnormal{(ii)}]
Show that an invertible fractional ideal is a projective~\(A\)-module.
\end{examsubparts}
\item
Suppose~\(A\) is a ring (not necessarily commutative, or with identity) and \(J\subseteq A\) is its Jacobson radical (the intersection of all regular maximal left (or right) ideals). Show that if \(J\ne A\) then \(A/J\) is semisimple.
\item
Suppose~\(K\) is a field and~\(A\) is a noncommutative semisimple~\(K\)-algebra of degree \(16\) (i.e. \(\dim_K(A)=16\)). List representatives of every possible isomorphism class of~\(A\) if
\begin{examsubparts}
\item[\textnormal{(i)}]
\(K=\mathbb{C}\),
\item[\textnormal{(ii)}]
\(K=\mathbb{R}\), or
\item[\textnormal{(iii)}]
\(K\) is a finite field.
\end{examsubparts}
\end{examproblems}
\end{document}
